\begin{graphicspathcontext}{{./chapters/ai/imgs/},{./chapters/ai/imgs/auto/},\old}

\adaptAIdefinitioncontent
\talkadaptivesidecite
\begin{frame}[c]{What is AI?}
	\talkadaptivesize
	\ifMinskyAI{
		\begin{definition}[Minsky 1961]
			AI is the science of \Emph{making machines do things} that would \Emph{require intelligence if done by men}
		\end{definition}
	}\fi
	\talkadaptivevspacea
	\ifOECDAI{
		\begin{definition}[OECD 2024]
			An AI is a machine-based system that, for \Emph{explicit or implicit goals}, deduces from the \Emph{inputs} it receives how to generate results such as \Emph{forecasts, content, recommendations or decisions} that can influence physical or virtual environments. \Emph{Different AI systems} vary in their levels of autonomy and adaptability once deployed
		\end{definition}
	}\fi
	\talkadaptivevspaceb
	\ifEUAIAct{
		\begin{definition}[EU IA Act]
			\Emph{AI system} means a machine-based system that is designed to operate with varying levels of autonomy and that may exhibit adaptiveness after deployment, and that, for \Emph{explicit or implicit objectives}, infers, from the \Emph{input} it receives, how to generate outputs such as \Emph{predictions, content, recommendations, or decisions} that can influence physical or virtual environments
		\end{definition}
	}\fi
\end{frame}

\end{graphicspathcontext}

\endinput

